\documentclass[10pt]{article}

\usepackage[utf8]{inputenc}

\usepackage[T1]{fontenc}

\usepackage{amsmath}

\usepackage{amsfonts}

\usepackage{amssymb}

\usepackage[version=4]{mhchem}

\usepackage{stmaryrd}



\begin{document}

Аналогично даются определения замкнутости, полноты и условия замкнутости для более общих пространств и мер.



Одним из важнейших вопросов теории О.р. является вопрос однозначного определения функции по ее коәфффициентам Фурье. Для пространств $L^{2}$ он самым тесным образом связан с выполнением равенства (2) для всех функций $f \in L^{2}$.



Для случая тригонометрич: сұстемы равенство (2) в 1805 было приведено (фактически без доказательства) М. Парсевалем (M. Parseval), а в 1828 Ф. Бессель установил, что



\[

\sum_{n=0}^{\infty} a_{n}^{2}(f) \leqslant \int_{a}^{b} f^{2}(x) d x

\]



ве равенство Бес селя). В 1896 А. М. Ляпунов доказал равенство (2) для интегрируемых по Риману функций, а потом П. Фату (P. Fatou) для случая $f \in L^{2}$.



В. А. Стекловым (1898-190') был поставлен вопрос о замкнутости общих онс и положительно решен для многих ортогональных систем (сферич. функции, собственные функции оператора Штурма - Лиувилля, системы ортогональных многочленов Әрмита, Лагерра, функции Ламе и др.).



Что касается неравенства (3), то оно оказалосг, справедливым для произвольных онс п функций $f \in L^{2}$.



В 1907 Ф. Рисс (F. Riesz) и Э. Фишер (E. Fischer) доказали, что для любой онс $\left\{\varphi_{n}\right\}$ и любой последовательности чисел $\left\{a_{n}\right\} \in l_{2}$ найдется функдия $f \in L^{2}$, для к-рой $a_{n}(f) \equiv\left(f, \varphi_{n}\right)=a_{n}$ и выполнено равенство (2). Из этой теоремы и неравенства Бесселя вытекает, что для любых онс полнота и замкнутость әквивалентны в пространстве $L^{2}$; замкнутость в пространствах $L^{p}$ с $1<p<\infty$ эквивалентна полноте в пространстве $L^{p^{\prime}}$, где $1 / p^{\prime}+1 / p=1$ (C. Банах, S. Banach, 1931).



Неравенство Бесселя и теорема Рисса - Фишера были распространены Г. Харди (G. Hardy), Дж. Литлвудом (J. Littlewood) и P. Пәли (R. Paley) на пространства $L^{p}$. Именно, пусть $\left\{\varphi_{n}\right\}-$ онс, $\left|\varphi_{n}(x)\right| \leqslant M$ и $1<p<2$. Тогда:



a) если $f \in L^{p}$, то



\[

\sum_{n=1}^{\infty}\left|\left(f, \varphi_{n}\right)\right|^{p} n^{p-2} \leqslant A\|f\|_{p}^{p}

\]



б) если дана последовательность $\left\{a_{n}\right\}$ с



\[

I \equiv \sum_{n=1}^{\infty}\left|a_{n}\right|^{p^{\prime}} n^{p^{\prime}-2}<\infty

\]



то найдется функция $f \in L^{p^{\prime}}$, для к-рой $\left(f, \varphi_{n}\right)-a_{n}$ и $\|f\|_{p^{\prime}}^{p^{\prime}} \leqslant A I$, где $A$ зависит лишь от $p$ и $M$.



\begin{enumerate}

  \setcounter{enumi}{1}

  \item Другой крупной проблемой теории О. р. является вопрос разложения функции в ряд по простым функциям, сходящийся к ней по норме того или иного пространства. Система элементов $\left\{\varphi_{n}\right\}$ из $B$-пространства $E$ наз. ба з и с ом $(0 ө з$ у с л о в ны м б а з и с о м), если каждый элемент $f \in E$ единственным образом представляется в виде ряда

\end{enumerate}



\[

f=\sum_{n=0}^{\infty} a_{n} \varphi_{n}, \quad a_{n}=a_{n}(f),

\]



сходящегося (безусловно сходящегося) к $f$ по норме пространства $E$.



Если $\left\{\varphi_{n}\right\}$ - базис в $E$, то $a_{n}(f)$ являются линеїными непрерывными функционалами в пространстве $E$ и в случаө $E=L^{p}(0,1)$ с $1<p<\infty$ имеют вид



\[

a_{n}(f)=\int_{0}^{1} f(t) \psi_{n}(t) d t

\]



где $\left\{\psi_{n}\right\}$ - базис пространства $L^{p^{\prime}}(0,1),\left\{\varphi_{n}, \psi_{n}\right\}-$ бпортонормированная система (С. Банах). В частности, если $\varphi_{n} \equiv \psi_{n}$, то есть $\left\{\varphi_{n}\right\}-$ онс, то ортого- нальный базис в $L^{p}$ автоматически является базисом во всех пространствах $L^{r}$, где $r-$ любое число между $p$ и $p^{\prime}$.



Исследования по указанной проблеме ведутся в двух направления $\mathbf{x}:$



a) по заданной онс $\left\{\varphi_{n}\right\}$ находятся те пространства, в к-рых $\left\{\varphi_{n}\right\}$ является базисом;



б) для заданного пространства $E$ отыскиваются в нем базисы или ортогональные базисы.



В обоих случаях исследуется взаимосвязь свойств функции $f$ и ее разложения.



Что касается тригонометрич. системы, то она не является базисом пространства непрерывных функций $C$ (П. Дюбуа-Реймон, P. Du Bois Reymond, 1876), но является базисом в простравствах $L^{P}$ с $1<p<\infty$ (М. Рисс, M. Riesz, 1927). Результат П. Дюбуа-Реймона был распространен на любые ограниченные в совокупности онс.



Ортонормированная система многочленов Лежандра являөтся базисом в пространствах $L^{p}$ при $p \in(4 / 3,4)$ и не является таковой в остальных пространствах $L^{q}$ (1946-52, Х. Поллард, Н. Pollard, Дж. Нейман, J. Neumann, и В. Рудин, W. Rudin).



В 1910 была построена онс $\left\{\chi_{m}\right\}_{m=0}^{\infty}$ такая, что всякая непрерыввая функция $f \in C(0,1)$ единственным образом раскладывается в равномерно сходящийся ряд Фурье по этой системе (А. Хаар, А. Наar). Однако система Xаapa $\left\{\chi_{m}\right\}$ не является базисом пространства $C(0,1)$, т. к. функции $\chi_{m}$ разрывны при $m \geqslant 1$. Проинтегрировав систему $\left\{\chi_{m}\right\}, \Gamma$. Фабер (G. Faber, 1910) установил, что система



\[

\left\{f_{n}(t)\right\} \equiv\left\{1, \int_{0}^{t} \chi_{m}(x) d x\right\}

\]



является базисом в пространстве $C(0,1)$ и тем самым был найден первый базис в пространстве непрерывных функций. Этот результат Г. Фабера был переоткрыт Ю. Шаудером (J. Schauder, 1927), к-рый указал также класс базисов пространства $C(0,1)$ типа базиса $\left\{f_{n}\right\}$; в честь последнего и введен термин "базис ІІаудера", хотя более справедливо было бы называть его «базис Фабера - Шаудера».



Построенные Г. Фабером и Ю. Шаудером базисы не являются ортогональными. Первый ортонормированный базис $\left\{F_{n}\right\}_{\text {в пространстве } C}(0,1)$ был найден $\Phi$. Франклином (Ph. Franklin, 1928), к-рый проортогонализировал методом Шмидта систему Фабера - Шаудера $\left\{f_{n}\right\}$ и получил $\left\{F_{n}\right\}$. На этом пути (ортогонализация и интегрирование) был введен и изучен новый класс базисов. Все ортонормированные базисы пространства $C(0,1)$ автоматически являются базисами во всех пространствах $L^{p}$ с $1 \leqslant p \leqslant \infty$.



Система Xаара $\left\{\chi_{m}\right\}$ является безусловным базисом во всех пространствах $L^{p}$ с $1<p<\infty(1931-37$, Р. Пэли, Ю. Марцинкевич, J. Marcinkiewicz). Аналогичный результат имеет место и для системы $\left\{F_{n}\right\}$ Франклина.



В пространствах $C$ и $L$ вообце нет безусловных базисов. Точно также не существует нормированных и ограниченных в совокупности безусловных базисов в пространствах $L^{p}$ при $1<p<\infty$ и $p \neq 2$.



\begin{enumerate}

  \setcounter{enumi}{2}

  \item Большой цикл исследований проведен по проблеме сходимости почти всюду тригонометрических и ортогональных рядов.

\end{enumerate}



В 1911 Н. Н. Лузин построил первый пример почти всюду расходящегося тригонометрич. ряда, коэффициенты к-рого стремятся к нулю. Такого типа ряд Фурье был построен А. Н. Колмогоровым (1923). Результат Н. Н. Јузина был распространен на произвольные полные онс, а результат А. Н. Колмогорова обобщен на множества положительной меры для ограниченвых в совокупности онс.





\end{document}